\scanpage{26}
\margin{\uline{propriétés} \uline{caractéristiques}}
\[
\left\{\begin{array}{ll}
\text{(i)} & \text{Compatibilité aux chgts de base admissibles}\\[4pt]
\text{(ii)} & \text{Si }Y=V(\varphi),\quad \underline\gamma_{Y/X}^{(n)}=d(\beta_{Y/X}^{(n)})\text{ comme dessus.}\\[4pt]
\text{(iii)} & \underline\gamma_{Y.Y'/X}^{(n)}=\underline\gamma_{Y/X}^{(n)}\cdot\underline\gamma_{Y'/X}^{(n)}
\end{array}\right.
\]
si $Y.Y'$ ``sans dim. excédentaire''

\[\Downarrow\]
(iv) Compatibilité pour $n$ variable

( \uline{NB} La ``bonne'' classe se trouve dans
$\Gamma(Y,\underline H^{2d}_Y(\mathbf Z_\ell(d)))$, il y en a une pour tt $\ell$ premier aux car. résiduelles de $Y$ ).

\uline{Application à la construction de classes de coh. associées à certains cycles.}

Supposons $\underline H^i_Y((\mathbf Z/n\mathbf Z)_X)=0$ \deleted{\unsure{qt}} quand $i<2d$

d'où également
\[
\underline H^i_Y(\mu_{nX}^{\otimes d})=0\qquad\text{pour }i<2d
\]
C'est la situation du th. de pureté, vrai

(i) $X,Y$ lisses sur une base $S$

(ii) $Y,X$ réguliers, $X$ excellent, car. nulles

(iii) $d=1$, $X$ normal.
\mspage{25}
